%==============================================================================
% CAPÍTULO 28 — CONCLUSÃO: DA LEITURA À PRÁTICA
% PARTE V — Futuro, Ética e Regulação da IA em Odontologia
% ÚLTIMO CAPÍTULO DA OBRA
%==============================================================================
\chapter{Conclusão: Da Leitura à Prática}
\label{cap:conclusao-final}

\nivelzero

Chegamos ao final de uma jornada de 500 laudas, 28 capítulos e incontáveis
linhas de código. Mas, como toda boa conclusão, este capítulo não é um ponto
final -- é um ponto de partida.

%-----------------------------------------------------------------------------
% SEÇÃO 28.1
%-----------------------------------------------------------------------------
\section{Síntese dos 27 Capítulos: O Que Construímos Juntos}
\label{sec:sintese}

\nivelum

Permita-me, leitor, recapitular o caminho que percorremos. Cada etapa foi
projetada para construir sobre a anterior, em uma arquitetura pedagógica
deliberada:

\subsection{Parte I -- Fundamentos Históricos e Tecnológicos (Capítulos 1--5)}

\subsubsection{As Estacas Históricas e a Mão na Massa}

Começamos pelo chão. Nos Capítulos 1 e 2, fincamos as estacas históricas:
compreendemos de onde viemos -- os primeiros sistemas computacionais
odontológicos, as três ondas da informatização, a emergência da IA como
ferramenta diagnóstica. Revisamos a evidência científica disponível, aprendendo
a separar o joio do trigo em um campo onde o \textit{hype} frequentemente
precede a substância.

Nos Capítulos 3, 4 e 5, colocamos a mão na massa: configuramos o Google Colab,
aprendemos Python aplicado à Odontologia, dominamos a estatística essencial
para não sermos enganados por números mal interpretados, e exploramos os
algoritmos de \textit{machine learning} clássico -- regressão logística,
\textit{random forests}, SVMs -- que, em muitos problemas odontológicos,
oferecem excelente relação custo-benefício.

\subsection{Parte II -- Teoria e Modelos de Deep Learning (Capítulos 6--11)}

\subsubsection{Redes Neurais, Segmentação e Diagnóstico}

Com as bases estabelecidas, mergulhamos fundo. Os Capítulos 6 e 7 desvendaram
as redes neurais artificiais e as CNNs -- as arquiteturas que revolucionaram a
visão computacional e, por extensão, a radiologia odontológica. O Capítulo 8
nos deu a U-Net e a segmentação semântica -- a capacidade de isolar dentes,
lesões e estruturas anatômicas com precisão de pixel.

Os Capítulos 9, 10 e 11 aplicaram esse conhecimento a problemas odontológicos
concretos: radiologia panorâmica e TCFC, classificação de lesões bucais e --
o ápice da Parte II -- detecção de câncer oral e desordens potencialmente
malignas, onde a IA pode, literalmente, salvar vidas.

\subsection{Parte III -- Engenharia de Software Odontológico com SDD/TDD (Capítulos 12--17)}

\subsubsection{De Usuário a Construtor de Sistemas Confiáveis}

A Parte III foi o divisor de águas. Aqui deixamos de ser ``usuários de IA''
para nos tornarmos ``construtores de sistemas confiáveis''. O Capítulo 12
introduziu o SDD (\textit{Specification-Driven Development}) como metodologia
para especificar antes de implementar. O Capítulo 13 nos deu o TDD
(\textit{Test-Driven Development}) -- escrever testes antes do código, garantindo
que cada linha produzida tenha uma razão de existir verificável.

Os Capítulos 14, 15, 16 e 17 consolidaram o ferramental do engenheiro de
\textit{software} odontológico: projetos guiados completos, interpretabilidade
de modelos (Grad-CAM, SHAP, LIME), auditoria de código no GitHub e -- talvez o
capítulo mais importante para a credibilidade científica -- reprodutibilidade.
Um modelo que só funciona no seu computador não é ciência; é anedota.

\subsection{Parte IV -- Aplicações Avançadas em Especialidades (Capítulos 18--23)}

\subsubsection{Especialização por Especialidade Odontológica}

Na Parte IV, especializamos. A Periodontia (Capítulo 18) ganhou modelos para
classificação de doença periodontal, medição de perda óssea radiográfica e
predição de progressão. A Implantodontia (Capítulo 19) recebeu planejamento
cirúrgico assistido por IA e predição de osseointegração. A Ortodontia
(Capítulo 20) foi transformada por cefalometria automatizada, predição de
crescimento facial e \textit{clear aligners} otimizados por IA.

O Capítulo 21 -- Gêmeos Digitais -- foi, para muitos leitores, o ponto alto
da obra: a criação de réplicas virtuais dinâmicas da boca do paciente, que
evoluem com ele ao longo do tempo e permitem simular tratamentos antes de
executá-los. Os Capítulos 22 e 23 nos levaram às fronteiras da IA generativa:
modelos fundacionais odontológicos e sistemas de RAG (\textit{Retrieval-Augmented
Generation}) para responder perguntas clínicas com evidências rastreáveis.

\subsection{Parte V -- Futuro, Ética e Regulação (Capítulos 24--28)}

\subsubsection{Do Código para a Sociedade}

A última parte -- que este capítulo encerra -- olhou para fora do código e para
dentro da sociedade. Aprendemos a publicar (Capítulo 24), a nos proteger
legalmente (Capítulo 25), a ensinar (Capítulo 26) e a vislumbrar o futuro
(Capítulo 27). A mensagem central da Parte V é inequívoca: \textbf{a excelência
técnica sem responsabilidade ética é uma ameaça; a responsabilidade ética sem
excelência técnica é impotência.}

%-----------------------------------------------------------------------------
% SEÇÃO 28.2
%-----------------------------------------------------------------------------

% === FIGURAS ADICIONADAS (OdontoIA Dark) ===
% ======================================================================
% FIGURA: Estrutura do Livro (Capítulo 28 — Conclusão)
% ======================================================================
\begin{figure}[H]
\centering
\begin{tikzpicture}[
    box/.style={rectangle, draw=blueaccent, fill=blueaccent!20, align=center, minimum width=3.5cm, align=center, minimum height=0.8cm, font=\small, rounded corners=3pt, align=center},
    arrow/.style={-latex, thick, blueaccent},
]
% Parte I
\node[align=center, box, fill=codegreen!25] (p1) at (0,0) {Parte I: Fundamentos\\\footnotesize 5 capítulos $\cdot$ $\bigstar$};
% Parte II
\node[align=center, box, fill=orangeaccent!25] (p2) at (0,-1.5) {Parte II: Teoria \& Modelos\\\footnotesize 6 capítulos $\cdot$ $\bigstar\bigstar\bigstar$};
% Parte III
\node[align=center, box, fill=red!20] (p3) at (0,-3.0) {Parte III: SDD/TDD\\\footnotesize 6 capítulos $\cdot$ $\bigstar\bigstar\bigstar\bigstar$};
% Parte IV
\node[align=center, box, fill=magenta!20] (p4) at (0,-4.5) {Parte IV: Aplicações\\\footnotesize 6 capítulos $\cdot$ $\bigstar\bigstar\bigstar\bigstar$};
% Parte V
\node[align=center, box, fill=cyanaccent!25] (p5) at (0,-6.0) {Parte V: Futuro \& Ética\\\footnotesize 5 capítulos $\cdot$ $\bigstar\bigstar\bigstar\bigstar\bigstar$};

% Setas de progressão
\draw[arrow] (p1) -- (p2);
\draw[arrow] (p2) -- (p3);
\draw[arrow] (p3) -- (p4);
\draw[arrow] (p4) -- (p5);

% Níveis à direita
\node[font=\small, anchor=west] at (4.5,0) {Nível 0-3: Leigo $\rightarrow$ Iniciante};
\node[font=\small, anchor=west] at (4.5,-1.5) {Nível 4-7: Intermediário clínico};
\node[font=\small, anchor=west] at (4.5,-3.0) {Nível 8-9: Avançado IA};
\node[font=\small, anchor=west] at (4.5,-4.5) {Nível 10+: Especialista};
\node[font=\small, anchor=west] at (4.5,-6.0) {PhD: Pesquisador};

% Total
\node[font=\bfseries, below=0.5cm of p5] {28 capítulos $\cdot$ 500+ laudas $\cdot$ 100+ artigos DOI};

\end{tikzpicture}
\caption{Estrutura progressiva do livro OdontoIA. O leitor avança do nível zero (leigo) ao nível PhD/pesquisador, passando por 28 capítulos organizados em 5 partes com dificuldade crescente.}
\label{fig:book-structure}
\end{figure}


\section{A Linhagem de Progressão do Leitor: De Onde Veio, a Onde Foi}
\label{sec:linhagem}

\niveldois

\subsubsection{Cinco Estágios de Transformação}

Se você seguiu este livro na ordem proposta, com disciplina e prática, sua
trajetória pode ser descrita em cinco estágios de transformação:

\begin{table}[H]
  \centering
  \caption{Linhagem de progressão do leitor do OdontoIA.}
  \label{tab:linhagem-leitor}
  \begin{tabularx}{\textwidth}{clX}
    \toprule
    \textbf{Estágio} & \textbf{Parte(s)} & \textbf{Você era... $\rightarrow$ Você se tornou...} \\
    \midrule
    1. Despertar & Parte I & Leigo curioso $\rightarrow$ Alfabetizado em IA
      odontológica \\
    2. Fundação & Parte II & Usuário passivo $\rightarrow$ Praticante de
      \textit{deep learning} \\
    3. Maestria & Parte III & Praticante $\rightarrow$ Engenheiro de software
      odontológico confiável \\
    4. Especialização & Parte IV & Generalista $\rightarrow$ Especialista em
      aplicações clínicas de IA \\
    5. Liderança & Parte V & Especialista $\rightarrow$ Líder ético, educador
      e visionário \\
    \bottomrule
  \end{tabularx}
\end{table}

Se você chegou até aqui, você não é mais o mesmo profissional que abriu este
livro pela primeira vez. Você agora:

\begin{itemize}
  \item Compreende como uma CNN processa uma radiografia panorâmica.
  \item Sabe especificar um sistema de IA odontológica antes de implementá-lo
    (SDD).
  \item Escreve testes antes do código (TDD) e não confia em nenhum modelo que
    não passe por todos eles.
  \item Conhece os \textit{checklists} CLAIM, PRISMA e TRIPOD e sabe como
    reportar sua pesquisa para periódicos Qualis A1.
  \item Compreende a LGPD, o PL 2338/2023 e a responsabilidade civil que
    acompanha o uso de IA na clínica.
  \item Sabe que a IA não substitui o dentista -- mas o dentista que ignora a
    IA será substituído pelo dentista que a domina.
  \item Tem um ecossistema -- o OdontoIA -- como parceiro nessa jornada.
\end{itemize}

\subsubsection{O Que Você Agora Sabe}

%-----------------------------------------------------------------------------
% SEÇÃO 28.3
%-----------------------------------------------------------------------------
\section{A Tríade: Evidência + Código + Ética}
\label{sec:triade}

\nivelcinco

Se este livro pudesse ser resumido em uma única frase, seria esta:

\destaque{Não basta que a IA funcione. É preciso que a IA funcione com
evidência, com código aberto e reprodutível, e com responsabilidade ética.}

Esta tríade -- \textbf{Evidência + Código + Ética} -- é o fio condutor que
atravessa todas as 500 laudas desta obra:

\begin{description}
  \item[Evidência:] Cada modelo que você construir deve estar ancorado na
    literatura odontológica. Não importa quão alta seja sua acurácia -- se o
    problema clínico que você resolveu não é relevante, ou se seu \textit{
    ground truth} é frágil, seu modelo é cientificamente vazio. Os Capítulos 2
    e 24 ensinaram como buscar, avaliar e produzir evidência de qualidade.

  \item[Código:] Código fechado não é ciência -- é prestidigitação. A
    reprodutibilidade é o alicerce do método científico, e \textit{software}
    proprietário a viola. Os Capítulos 12--17 forneceram a metodologia (SDD/TDD)
    e as ferramentas (GitHub, Colab, Zenodo) para que cada modelo que você
    construir seja verificável, reproduzível e auditável por qualquer pessoa,
    em qualquer lugar do mundo.

  \item[Ética:] A IA odontológica lida com os dados mais íntimos de um ser
    humano -- seu sorriso, sua face, sua saúde. Violar a privacidade, perpetuar
    vieses, delegar decisões que deveriam ser humanas -- esses não são ``danos
    colaterais'' da inovação; são falhas de projeto. O Capítulo 25 forneceu o
    arcabouço legal e ético para que sua prática seja não apenas eficaz, mas
    justa.
\end{description}

\subsubsection{Citação de Apoio}

\citacaoativa{doi-3389-fdmed-2023-1085251}{10.3389/fdmed.2023.1085251}{%
  ``A tríade evidência-código-ética não é uma restrição à inovação -- é a
  condição para que a inovação seja durável, escalável e digna da confiança
  que pacientes e profissionais depositam na Odontologia.''}

%-----------------------------------------------------------------------------
% SEÇÃO 28.4
%-----------------------------------------------------------------------------
\section{Call to Action: O Que Fazer Depois de Ler Este Livro}
\label{sec:call-to-action}

\niveldois

Você concluiu a leitura. O conhecimento está em sua mente. Agora, a pergunta
que realmente importa: \textbf{o que você vai fazer com ele?}

Eis um roteiro de ação, organizado por perfil. Escolha o seu -- ou crie o seu
próprio:

\subsection{Se Você é Estudante de Odontologia...}

\begin{enumerate}
  \item \textbf{Crie seu primeiro notebook ainda esta semana.} Vá ao Google
    Colab, abra um notebook em branco e execute \texttt{import tensorflow as tf;
    print(tf.\_\_version\_\_)}. A primeira linha de código é a mais difícil; as
    próximas virão naturalmente.

  \item \textbf{Convença seu professor a integrar IA à disciplina.} Leve este
    livro para a sala de aula. Mostre o Capítulo 26. Ofereça-se para apresentar
    um seminário sobre IA aplicada à especialidade da disciplina. Seja o
    catalisador.

  \item \textbf{Participe da comunidade OdontoIA.} Entre no Discord, no grupo
    de WhatsApp. Pergunte, responda, compartilhe. A comunidade existe para
    apoiar sua jornada.

  \item \textbf{Pense no seu TCC.} Um Trabalho de Conclusão de Curso envolvendo
    IA odontológica -- com metodologia SDD/TDD, código no GitHub, dataset
    documentado, checklist CLAIM -- não é apenas um requisito curricular; é um
    artigo publicável e um diferencial de currículo.
\end{enumerate}

\subsection{Se Você é Cirurgião-Dentista Clínico...}

\begin{enumerate}
  \item \textbf{Implemente uma ferramenta de IA no seu consultório nos próximos
    3 meses.} Pode ser um sistema de detecção de cáries em radiografias, um
    \textit{software} de planejamento ortodôntico com IA, ou simplesmente o
    Odontinho® como ferramenta de educação de pacientes. Dê o primeiro passo.

  \item \textbf{Audite a conformidade LGPD do seu consultório.} Use o template
    do Capítulo 25. Verifique seus termos de consentimento. Documente o uso de
    IA no prontuário.

  \item \textbf{Torne-se um produtor de dados de qualidade.} Comece a anotar
    suas radiografias e imagens clínicas de forma estruturada. Participe de
    estudos multicêntricos. A comunidade de IA odontológica precisa de dados
    anotados por especialistas brasileiros -- e você pode fornecê-los.

  \item \textbf{Seja um multiplicador.} Conte para seus colegas o que aprendeu.
    Indique este livro. Apresente casos clínicos em congressos mostrando como
    a IA auxiliou seu diagnóstico. Combata o ceticismo com evidência e o
    \textit{hype} com sobriedade.
\end{enumerate}

\subsection{Se Você é Pesquisador ou Pós-Graduando...}

\begin{enumerate}
  \item \textbf{Publique um artigo seguindo os padrões deste livro.} CLAIM 2024,
    código no GitHub com DOI (Zenodo), \textit{dataset} documentado com
    \textit{Datasheets for Datasets}, análise de vieses por subgrupo. Eleve o
    padrão da pesquisa brasileira em IA odontológica.

  \item \textbf{Colabore com o OdontoIA.} Submeta \textit{pull requests} no
    GitHub, contribua com \textit{datasets}, escreva tutoriais, grave episódios
    para o Podcast do Odontinho. A plataforma é aberta e sua contribuição será
    creditada.

  \item \textbf{Oriente alunos de iniciação científica em IA odontológica.}
    Use este livro como roteiro. Cada capítulo é uma porta de entrada para um
    projeto de pesquisa.

  \item \textbf{Candidate-se a editais de fomento.} CNPq, CAPES, FAPEAL,
    FAPESP -- há linhas de financiamento para IA em saúde. Este livro e o
    ecossistema OdontoIA fornecem a fundamentação e a infraestrutura para
    projetos competitivos.
\end{enumerate}

\subsection{Se Você é Desenvolvedor ou Cientista da Computação...}

\begin{enumerate}
  \item \textbf{Mergulhe na Odontologia.} Leia os Capítulos 1 e 2 com atenção.
    Entenda a anatomia dental, as patologias orais, o fluxo clínico. O melhor
    modelo de IA para Odontologia será construído por quem entende de ambos os
    mundos.

  \item \textbf{Construa pontes.} Procure um dentista para colaborar. Ofereça
    sua expertise técnica em troca de expertise clínica. As melhores inovações
    nascem na interseção.

  \item \textbf{Contribua com código aberto.} O ecossistema OdontoIA precisa
    de \textit{front-end}, \textit{back-end}, \textit{MLOps}, documentação,
    testes, segurança. Há espaço para todas as habilidades.
\end{enumerate}

%-----------------------------------------------------------------------------
% SEÇÃO 28.5
%-----------------------------------------------------------------------------
\section{Comunidade OdontoIA: Como Continuar Aprendendo}
\label{sec:comunidade}

\nivelzero

O aprendizado não termina na última página deste livro. O ecossistema OdontoIA
oferece múltiplos caminhos para continuar sua jornada:

\begin{itemize}
  \item \textbf{Portal OdontoIA (\url{https://odontoia.com.br}):} Artigos
    semanais, tutoriais, atualizações de literatura e novos módulos do
    Odontinho®.

  \item \textbf{GitHub (\url{https://github.com/odontoia}):} Código-fonte de
    todos os notebooks, modelos e ferramentas. \textit{Issues}, \textit{pull
    requests} e discussões técnicas.

  \item \textbf{Discord e WhatsApp:} Comunidade ativa de aprendizagem
    colaborativa, sessões semanais de \textit{coding}, revisão por pares,
    desafios mensais.

  \item \textbf{Podcast do Odontinho (Spotify):} Episódios quinzenais com
    entrevistas, revisões de literatura e discussões sobre IA odontológica.

  \item \textbf{Odontinho Games:} Plataforma gratuita de gamificação para
    treinamento contínuo em diagnóstico, farmacologia e planejamento.

  \item \textbf{YouTube OdontoIA:} Tutoriais em vídeo, \textit{webinars} ao
    vivo e gravações de cursos e palestras.

  \item \textbf{Newsletter mensal:} Resumo dos principais artigos de IA
    odontológica publicados no mês, com curadoria e comentários.
\end{itemize}

\destaque{O OdontoIA não é um produto -- é uma comunidade. E você, leitor, está
convidado a participar, contribuir e liderar.}  

%-----------------------------------------------------------------------------
% SEÇÃO 28.6
%-----------------------------------------------------------------------------
\section{Carta Final do Autor Edson Laranjeiras ao Leitor}
\label{sec:carta-autor}

\nivelzero

\textit{Caro leitor, cara leitora,}

\textit{Quando comecei a escrever este livro, eu tinha uma ideia clara do que
queria entregar: uma obra que fosse, ao mesmo tempo, rigorosa e acolhedora;
técnica e inspiradora; que partisse do zero absoluto e chegasse às fronteiras
do conhecimento.}

\textit{O que eu não sabia -- e só descobri ao longo das mais de 500 laudas que
você acaba de percorrer -- era que este livro seria, também, uma carta de amor
à Odontologia brasileira.}

\textit{Amor porque este livro é gratuito. Amor porque cada linha de código está
em repositório aberto. Amor porque cada capítulo foi escrito pensando no
estudante que não tem dinheiro para um simulador de USD 100.000, no dentista
do interior que atende 40 pacientes por dia no SUS, na pesquisadora que luta
para publicar com os recursos limitados de uma universidade pública.}

\textit{Este livro é minha tentativa de retribuir um pouco do que a Odontologia
me deu. A profissão que me formou, os pacientes que me ensinaram, os colegas
que me inspiraram, os alunos que me desafiaram -- todos estão nestas páginas.}

\textit{Se você chegou até aqui, quero lhe fazer um pedido: não deixe este
conhecimento parado. Use-o. Aplique-o. Critique-o. Melhore-o. Ensine-o.
O Brasil precisa de dentistas que compreendam IA. O SUS precisa de sistemas
inteligentes que reduzam desigualdades. A ciência odontológica precisa de
pesquisadores que publiquem com reprodutibilidade e ética.}

\textit{Você é, agora, parte dessa vanguarda.}

\textit{Obrigado pela confiança de me acompanhar até aqui. Nos encontramos na
comunidade OdontoIA, nos congressos, nos artigos, nas trincheiras do SUS e nas
salas de aula. A jornada continua -- e mal começou.}

\begin{flushright}
\textit{Com gratidão e esperança,}

\textit{Edson Laranjeiras} \\
Cirurgião-Dentista | Criador do OdontoIA e do Odontinho® \\
Maceió, Alagoas, Brasil \\
Julho de 2026
\end{flushright}

%-----------------------------------------------------------------------------
% SEÇÃO 28.7
%-----------------------------------------------------------------------------
\section{Dedicatória ao Futuro da Odontologia Brasileira}
\label{sec:dedicatoria-futuro}

\nivelzero

\begin{center}
\itshape

\bigskip

Este livro é dedicado:

\medskip

Ao estudante de Odontologia que, nesta noite, \\
abrirá o Google Colab pela primeira vez. \\
Que seu primeiro \texttt{import} seja o início de uma revolução silenciosa.

\medskip

Ao cirurgião-dentista do SUS que atende 40 pacientes por dia \\
e ainda encontra forças para se atualizar. \\
Que a IA seja sua aliada, não sua substituta.

\medskip

À pesquisadora brasileira que luta para publicar \\
com equipamentos limitados e financiamento escasso. \\
Que este livro lhe dê método, rigor e esperança.

\medskip

Ao paciente do sistema público que espera meses por uma consulta \\
e merece o melhor que a tecnologia pode oferecer. \\
Que a IA reduza o tempo entre o sintoma e o tratamento.

\medskip

Aos meus pais, que me ensinaram que o conhecimento \\
só tem valor quando compartilhado.

\medskip

E a você, leitor, que dedicou seu tempo e sua atenção a esta obra. \\
Que ela tenha valido cada minuto.

\bigskip

\textit{A Odontologia brasileira merece o futuro que estamos construindo.} \\
\textit{E esse futuro começa agora.}

\end{center}

%-----------------------------------------------------------------------------
% SEÇÃO 28.8
%-----------------------------------------------------------------------------
\section{``Se Este Livro Inspirou ao Menos Um Dentista a Abrir o Colab...''}
\label{sec:fecho}

\nivelzero

No Prefácio, prometi que este livro seria diferente. Prometi que não seria mais
um livro de Odontologia que menciona IA superficialmente, nem mais um livro de
IA que ignora a realidade clínica. Prometi código real, dados reais, ética real.

Se, ao final destas 500 laudas, você sente que essa promessa foi cumprida, então
este livro cumpriu sua missão.

Mas a verdadeira medida de sucesso desta obra não está nestas páginas -- está
no que acontecerá depois delas. Está no notebook do Google Colab que você abrirá
amanhã. Está na primeira radiografia que você submeterá a um modelo treinado
por você mesmo. Está no artigo que você publicará seguindo o CLAIM 2024. Está
no paciente que receberá um diagnóstico mais precoce porque você decidiu, hoje,
aprender IA.

\textbf{Se este livro inspirou ao menos um dentista a abrir o Colab e executar
seu primeiro \texttt{import tensorflow as tf}, sua missão terá sido cumprida.}

Se inspirou mais do que um -- e tenho esperança de que inspirou --, então
estamos diante do início de uma transformação geracional na Odontologia
brasileira.

\medskip

\begin{center}
\Large

A Odontologia mudou.

A IA chegou.

E você está pronto.

\bigskip

\textbf{Fim da Parte V -- e do Livro.}

\medskip

Mas o começo da sua jornada.

\vspace{1cm}

\includegraphics[width=0.3\textwidth]{imagens/odontoia-logo.png}

\url{https://odontoia.com.br}

\bigskip

\textit{``Odontologia + Inteligência Artificial: o futuro começa agora.''}

\vspace{0.5cm}

\today
\end{center}

%-----------------------------------------------------------------------------
% SEÇÃO 28.9
%-----------------------------------------------------------------------------
\section{Última Síntese: O Que Levamos Deste Livro}
\label{sec:ultima-sintese}

\niveldois

\begin{itemize}
  \item A Odontologia e a Inteligência Artificial não são campos separados --
    são aliados naturais na busca por diagnósticos mais precisos, tratamentos
    mais eficazes e cuidados mais acessíveis.

  \item \textbf{Evidência} (Capítulos 2, 24), \textbf{Código} (Capítulos 3--17)
    e \textbf{Ética} (Capítulo 25) formam a tríade indissociável que deve guiar
    toda aplicação de IA em saúde bucal.

  \item O \textbf{OdontoIA} e o \textbf{Odontinho®} não são apenas ferramentas
    -- são a materialização de uma visão: a de que o conhecimento de ponta deve
    ser gratuito, aberto e acessível a todos os profissionais de Odontologia,
    independentemente de sua localização geográfica ou condição socioeconômica.

  \item O \textbf{futuro} (Capítulo 27) não é algo a temer, mas a construir.
    Bioimpressão 3D, nanotecnologia, IoT, teleodontologia e Odontologia 5P são
    horizontes alcançáveis -- e o Brasil pode ser protagonista.

  \item A \textbf{educação} (Capítulo 26) é a chave. Formar dentistas
    alfabetizados em IA, críticos, éticos e criativos é a missão central deste
    livro e do ecossistema OdontoIA.

  \item E, acima de tudo, \textbf{você, leitor, é o recurso mais valioso desta
    obra.} Seu tempo, sua atenção, sua curiosidade e sua coragem de aprender
    algo novo são o que dará vida a estas páginas.
\end{itemize}

\subsubsection{Mensagem Final}

\vspace{1cm}

\begin{center}
\begin{mdframed}[
  linecolor=corPrimaria,
  linewidth=3pt,
  backgroundcolor=corPrimaria!5,
  roundcorner=4pt,
  innertopmargin=15pt,
  innerbottommargin=15pt
]
  \Large\textbf{\textcolor{corPrimaria}{Obrigado.}}

  \bigskip

  \normalsize
  Este livro é para você. \\
  Este livro é por você. \\
  Este livro é com você.

  \bigskip

  \textit{A jornada continua.} \\
  \textit{-- Edson Laranjeiras e toda a comunidade OdontoIA}
\end{mdframed}
\end{center}

\endinput
